% !TeX root = ../presentacion.tex
%=======================================================================
%  config/plantilla.tex — mecanismos propios de la plantilla
%
%  Dos interruptores gobiernan qué se ve. ANTES DE EXPONER pon los dos en
%  «false»: desaparecen las cajas de guion y los marcadores quedan en negro,
%  sin que haya que borrar nada del texto.
%=======================================================================

%------------------------ Interruptores --------------------------------
%  Estado actual: los dos en «false». La presentación está cerrada y lista
%  para proyectar. Ponlos en «true» solo si se vuelve a maquetar.
\newif\ifguia          \guiafalse       % \guiatrue  muestra los guiones al maquetar
\newif\ifmarcadores    \marcadoresfalse % \marcadorestrue resalta en ámbar lo que falte

%--------------------------- Guion de la lámina ------------------------
% \guion{...} — qué se cuenta de viva voz en esta lámina. Se apaga con
% \guiafalse. Es distinto de \note{...}: el guion se IMPRIME en la lámina
% mientras se prepara la charla (para verlo al maquetar); la nota va a la
% pantalla del ponente y nunca al proyector.
\newcommand{\guion}[1]{%
  \ifguia
    \begin{tcolorbox}[enhanced,
        colback=guiaback,colframe=guiaframe,boxrule=0.4pt,arc=1pt,
        left=5pt,right=5pt,top=3pt,bottom=3pt,
        fontupper=\scriptsize\itshape]
      \textbf{\textup{Guion.}}\enspace #1
    \end{tcolorbox}%
  \fi}

% \ph{...} — marcador de relleno, el equivalente de «[Insertar X]».
% En ámbar mientras se prepara; en negro al exponer, para que salte a la
% vista si alguno se quedó sin rellenar.
\newcommand{\ph}[1]{%
  \ifmarcadores{\color{marcador}\textsf{[#1]}}\else\textsf{[#1]}\fi}

%------------------------- Cabecera de sección -------------------------
% \bloque{Título}{minutos}{quién expone}{qué criterio de la rúbrica defiende}
%
% Genera la lámina separadora de metropolis y, debajo del título, la ficha
% del bloque. El minutaje sale impreso a propósito: obliga a que la suma de
% los bloques quepa en el tiempo asignado, en vez de descubrirlo ensayando.
\newcommand{\bloque}[4]{%
  \section{#1}%
  \ifguia
    \begin{frame}[plain]{#1}
      \begin{tcolorbox}[enhanced,
          colback=acentoClaro,colframe=acento,boxrule=0.6pt,arc=2pt,
          left=8pt,right=8pt,top=6pt,bottom=6pt]
        \small
        \textbf{Duración:} #2 \quad\textbullet\quad
        \textbf{Expone:} #3\\[2pt]
        \textbf{Criterio de la rúbrica:} #4
      \end{tcolorbox}
    \end{frame}%
  \fi}

%----------------------- Identificadores trazables ---------------------
% Los mismos que el informe, para que el público que tenga el E1 delante
% reconozca el identificador sin traducirlo.
%
%   \rn{02}      RN-02   regla de negocio
%   \adrref{01}  ADR-01  decisión de arquitectura
%   \vista{02}   V-02    vista arquitectónica
%   \crit{3}     CA-3    criterio de aceptación
%   \ent{1}      E1      entregable del proyecto
\newcommand{\rn}[1]{\textbf{RN-#1}}
\newcommand{\adrref}[1]{\textsf{ADR-#1}}
\newcommand{\vista}[1]{\textsf{V-#1}}
\newcommand{\crit}[1]{\textsf{CA-#1}}
\newcommand{\ent}[1]{\textbf{E#1}}

% Identificador técnico en monoespaciado: \id{ms-reservas}
% \smaller y no \small: el identificador tiene que encoger respecto al texto
% que lo rodea, sea este \normalsize o \scriptsize. Con un tamaño fijo, un
% \id dentro de una lista pequeña se veía MÁS grande que la lista.
\newcommand{\id}[1]{\texttt{\smaller #1}}

%-------------------------- Marcas de estado ---------------------------
% Para las tablas de cumplimiento. Nunca marques \si{} algo que no hayas
% visto correr: un criterio dado por bueno que falla en la demo cuesta más
% que uno declarado pendiente.
\newcommand{\si}{{\color{ok}\ding{51}}}
\newcommand{\no}{{\color{alerta}\ding{55}}}
\newcommand{\parcial}{{\color{marcador}\ding{108}}}

%------------------------- Figura a lámina completa --------------------
% \laminafigura{archivo}{Pie de figura}
%
% Un diagrama ocupa toda la lámina, sin título ni pie de página que le
% roben altura. Es como se proyectan los C4: si hay que entrecerrar los
% ojos para leer una etiqueta, la lámina no sirve.
\newcommand{\laminafigura}[2]{%
  \begin{frame}[plain]
    \begin{center}
      \includegraphics[width=\textwidth,height=0.84\textheight,keepaspectratio]{#1}\\[2pt]
      {\scriptsize\color{textoSuave} #2}
    \end{center}
  \end{frame}}

% \laminafiguratitulo{Título}{archivo}{Pie}
% Igual, pero conservando el título: útil cuando el diagrama necesita que se
% diga en la propia lámina qué se está mirando.
\newcommand{\laminafiguratitulo}[3]{%
  \begin{frame}{#1}
    \begin{center}
      \includegraphics[width=\textwidth,height=0.66\textheight,keepaspectratio]{#2}\\[2pt]
      {\scriptsize\color{textoSuave} #3}
    \end{center}
  \end{frame}}

%--------------------------- Cifra destacada ---------------------------
% \cifra{11}{contenedores} — un dato grande y su rótulo. Para las láminas de
% resultados, donde un número se recuerda y una frase no.
\newcommand{\cifra}[2]{%
  \begin{minipage}[t]{\linewidth}
    \centering
    {\fontsize{34}{38}\selectfont\bfseries\color{acento} #1}\\[2pt]
    {\footnotesize\color{textoSuave} #2}
  \end{minipage}}

%------------------------- Bloque de decisión (ADR) --------------------
% \decisionslide{01}{Título}{Decisión}{Alternativa descartada}
% Versión de lámina del ADR del informe: en pantalla solo caben la decisión
% y lo que se descartó. El contexto y las consecuencias, en el E1.
\newcommand{\decisionslide}[4]{%
  \begin{block}{\adrref{#1} — #2}
    \scriptsize
    \textbf{Decisión.} #3\\[2pt]
    \textbf{Descartado.} #4
  \end{block}}

%----------------------- Mapa de arquitectura --------------------------
% \caja{tinta}{Rótulo}{Segunda línea} — un recuadro del mapa de la lámina
% de arquitectura. Existe para que las cinco filas del mapa se escriban
% igual y queden alineadas: una tabla de \colorbox sueltos se descuadra en
% cuanto una celda lleva dos líneas y otra una.
\newcommand{\caja}[3]{%
  \setlength{\fboxsep}{3pt}%
  \colorbox{#1}{\parbox[c][19pt][c]{\dimexpr\linewidth-6pt\relax}{%
    \centering\scriptsize\textbf{#2}%
    \ifx\relax#3\relax\else\\[-2pt]{\tiny\color{textoSuave}#3}\fi}}}

% Flecha de una fila a la siguiente, con su rótulo al lado.
\newcommand{\bajada}[1]{%
  \makebox[\linewidth]{\scriptsize\color{textoSuave}$\downarrow$\enspace #1}}
